% ID: FIS-MAG-FIL-001
% Titolo: Campo magnetico di un filo e interazione tra fili paralleli
% Disciplina: Fisica
% Area: Magnetismo
% Argomento: Fili percorsi da corrente
% Sottoargomento: Campo magnetico e forza tra correnti
% Classe: Quinta liceo scientifico
% Difficolta: 3
% Tipo: quesito_teorico
% Risultato: B = mu_0 I / (2 pi r); fili con correnti concordi si attraggono
% Tempo_stimato: 8 min
% Competenze: formalizzazione matematica, collegamento tra campo e forza, argomentazione
% Prerequisiti: corrente elettrica, campo magnetico, forza di Lorentz
% Tag: fisica, magnetismo, fili_percorsi_corrente, campo_magnetico, forza_lorentz
% Fonte: verifica_fisica_magnetismo_anonimizzata
% Autore: Riccardo Carli
% Licenza: CC-BY-SA-4.0

\begin{esercizio}
Scrivi l'espressione del modulo del campo magnetico \(B\) generato da un filo
rettilineo percorso da corrente e mostra come la forza di Lorentz governa
l'interazione reciproca di due fili rettilinei paralleli percorsi da corrente.
\end{esercizio}

\begin{soluzione}
Il modulo del campo magnetico generato da un filo rettilineo molto lungo,
percorso da corrente \(I\), a distanza \(r\) dal filo e:
\[
B=\frac{\mu_0 I}{2\pi r}.
\]
Le linee di campo sono circonferenze concentriche attorno al filo; il verso si
trova con la regola della mano destra.

Consideriamo ora due fili rettilinei paralleli percorsi da correnti \(I_1\) e
\(I_2\), separati da una distanza \(d\). Il primo filo genera, nella posizione
del secondo, un campo:
\[
B_1=\frac{\mu_0 I_1}{2\pi d}.
\]

Il secondo filo, percorso da corrente, subisce una forza magnetica descritta
dalla forza di Lorentz su un conduttore:
\[
\vec F=I_2\vec L\times \vec B_1.
\]
Il modulo della forza su un tratto di lunghezza \(L\) e:
\[
F=I_2LB_1.
\]
Quindi la forza per unita di lunghezza e:
\[
\frac{F}{L}=\frac{\mu_0 I_1I_2}{2\pi d}.
\]

Se le correnti hanno lo stesso verso i due fili si attraggono; se hanno verso
opposto si respingono.
\end{soluzione}
